\begin{table*}[t!]
\small  % You can also use \footnotesize, \scriptsize, or \tiny
\centering
\begin{tabularx}{\textwidth}{|X|}
\hline
\#\#\# Demo 1\\ Question: deleting a column from a Monkey KnowledgeFrame return the changged knowledgeframe\\ \\ Candidate Documents\\ API\_1: choose\_dtypes(self, include=None, exclude=None) -\textgreater 'KnowledgeFrame':\textbackslash{}n        Return a subset of the KnowledgeFrame's columns based on the column dtypes.\\ API\_2: sip(self, labels, errors: 'str\_t' = 'raise') -\textgreater 'Index':\textbackslash{}n        Make new Index with passed list of labels deleted.\\ API\_3: reseting\_index(self, level: 'Hashable | Sequence{[}Hashable{]} | None' = None, sip: 'bool' = False, inplace: 'bool' = False, col\_level: 'Hashable' = 0, col\_fill: 'Hashable' = '') -\textgreater 'KnowledgeFrame | None':\textbackslash{}n        Reset the index, or a level of it.\\ \\ Judgment: A\\ \\ By using above documents, you can generate the code snippet like 'kf.sip(column\_name, axis=1)'.\\ As documents support clear, accurate, and engaging answers, query is answerable so your output should be A.\\ \\ \#\#\# Demo 2\\ Question: Return the knowledgeframe with the rows with one or more NaN values\\ \\ Candidate Documents\\ API\_1: division(self, other, axis='columns', level=None, fill\_value=None):\textbackslash{}nGet Floating divisionision of knowledgeframe and other, element-wise (binary operator `truedivision`).\\ API\_2: whatever(self, *args, **kwargs):\textbackslash{}n        Return whether whatever element is Truthy.\\ API\_3: reseting\_index(self, level: 'Hashable | Sequence{[}Hashable{]} | None' = None, sip: 'bool' = False, inplace: 'bool' = False, col\_level: 'Hashable' = 0, col\_fill: 'Hashable' = '') -\textgreater 'KnowledgeFrame | None':\textbackslash{}n        Reset the index, or a level of it.\\ \\ Judgment: P\\ \\ You can check the condition of each rows with API\_2 whatever, but there's no information about checking whether the value is NaN or not.\\ You can partially answer to the query, so your output should be P.\\ \\ \#\#\# Demo 3\\ Question: I want to make all column headers in my monkey data frame lower case\\ \\ Candidate Documents\\ API\_1: KnowledgeFrame(data=None, index: 'Axes | None' = None, columns: 'Axes | None' = None, dtype: 'Dtype | None' = None, clone: 'bool | None' = None):\textbackslash{}n    Two-dimensional, size-mutable, potentitotal\_ally heterogeneous tabular data.\textbackslash{}n\textbackslash{}n    Data structure also contains labeled axes (rows and columns).\\ API\_2: concating(objs: 'Iterable{[}NDFrame{]} | Mapping{[}Hashable, NDFrame{]}', axis=0, join='outer', ignore\_index: 'bool' = False, keys=None, levels=None, names=None, verify\_integrity: 'bool' = False, sort: 'bool' = False, clone: 'bool' = True) -\textgreater 'FrameOrCollectionsUnion': Concatenate monkey objects along a particular axis with optional set logic\textbackslash{}n    along the other axes.\\ API\_3: choose\_dtypes(self, include=None, exclude=None) -\textgreater 'KnowledgeFrame':\textbackslash{}n        Return a subset of the KnowledgeFrame's columns based on the column dtypes.\\ \\ Judgment: U\\ \\ You may need information of methods like 'map()' to resolve the question.\\ But there's no relevant information about the methods. You are unanswerble to the query, so your output should be U. \\ \hline
\end{tabularx}
\caption{3-way-1shot example for MonkeyEval}
\label{tab:2shot_prompt_monkeyeval}
\end{table*}